\chapter{Triển khai các Task}\label{ch:trienkhai-se}
Ảnh minh họa là bằng chứng thu thập trên bản triển khai thật ngày 23/10/2026 (\path{docs/evidence}), tài khoản khách \texttt{customer@brewlite.com}.

\section{Kho mã và quản lý tài liệu}\label{sec:se-repo}
Kho \texttt{tthan-prog/BrewLite} theo cấu trúc monorepo: \path{apps/web}, \path{apps/api}, \path{tests/e2e}, \path{docs}, \path{.github/workflows}, \path{docker-compose.yml}. Nhánh \texttt{main} được bảo vệ: phải qua PR, hai check \texttt{api} và \texttt{web} phải xanh, cần một approval, không force-push. Tài liệu Scrum, yêu cầu, sơ đồ và bằng chứng nằm trong \path{docs}.

\section{Hiện thực các chức năng chính}\label{sec:se-hienthuc}
\subsection{Task 1: Khởi tạo dự án}
Có README, \texttt{.gitignore}, \texttt{.env.example}, \texttt{docker-compose.yml} (ba dịch vụ db, backend, frontend) và endpoint \texttt{/health}.

\subsection{Task 2--4: Sản phẩm, menu, chi tiết}
\texttt{GET /products} và \texttt{GET /products/:id} đọc từ PostgreSQL (seed 12 sản phẩm, 4 topping); id không tồn tại trả 404, id sai định dạng trả 400. Trang menu gọi API và xử lý ba trạng thái loading, rỗng, lỗi (\cref{fig:menu}). Trang chi tiết cho chọn size, mức đường, mức đá và topping (\cref{fig:size}); giá do server tính nên khách không sửa được từ phía trình duyệt.
\hinh{ev01-menu.jpg}{.92\linewidth}{Thực đơn tải từ API}{fig:menu}
\hinh{ev02-chon-size.jpg}{.92\linewidth}{Chọn size, mức đường và mức đá}{fig:size}

\subsection{Task 5: Giỏ hàng}
Giỏ hàng thêm, sửa số lượng, xóa món, hiển thị tổng tiền và huy hiệu số lượng. Tổng tiền $=\sum$ (đơn giá + phụ phí size + topping) $\times$ số lượng (\cref{fig:cart}).
\hinh{ev03-gio-hang.jpg}{.92\linewidth}{Giỏ hàng: 10 ly, 34.000 đồng/ly, tạm tính 340.000 đồng}{fig:cart}

\subsection{Task 6--7: Tạo đơn và xác thực}
Đăng ký và đăng nhập bằng email, mật khẩu băm bcrypt, JWT có hạn (mặc định một ngày); trùng email trả 409, sai mật khẩu trả 401. \texttt{POST /orders} yêu cầu JWT, lấy \texttt{userId} từ token (trường \texttt{userId} trong body bị bỏ qua), kiểm tra dữ liệu, trừ tồn kho và tạo đơn PENDING trong một transaction, và bắt buộc có header \texttt{Idempotency-Key}.

\subsection{Task 8--9: Thanh toán, xác nhận, lịch sử}
\texttt{POST /payments} xử lý thanh toán Ví hoặc Thẻ mô phỏng: thành công đưa đơn sang PAID; thất bại đưa sang PAYMENT\_FAILED (API trả 200 kèm trạng thái đơn) và cho thử lại. Khách xem xác nhận đơn và lịch sử tại \texttt{/orders/history}.
\hinh{ev04-thanh-toan.jpg}{.92\linewidth}{Thanh toán bằng thẻ mô phỏng, mã \texttt{WELCOME10}}{fig:pay}
\hinh{ev05-xac-nhan-don1-paid.jpg}{.92\linewidth}{Đơn \#1 PAID, 306.000 đồng (340.000 đồng giảm 10\%)}{fig:paid}
\hinh{ev09-thanh-toan-that-bai.jpg}{.92\linewidth}{Giả lập thanh toán thất bại: đơn \#2 PAYMENT\_FAILED, có nút ``Thử thanh toán lại''}{fig:failed}
\hinh{ev10-thu-lai-don2-paid.jpg}{.92\linewidth}{Thử lại thành công: đơn \#2 PAID, 585.000 đồng}{fig:retry}
\hinh{ev06-lich-su-don1.jpg}{.92\linewidth}{Lịch sử đơn của khách}{fig:history}

\subsection{Task 10: Nghiệp vụ nâng cao}
\begin{description}[leftmargin=0pt,style=nextline]
\item[Máy trạng thái đơn.] Hàm chuyển trạng thái ở \path{apps/api/src/orders/order-state.ts} (\cref{tab:state}); unit test duyệt mọi cặp trạng thái, test tích hợp kiểm tra API trả 400 khi chuyển sai.
\item[Idempotency.] Cặp \texttt{(userId, idempotencyKey)} có ràng buộc duy nhất cùng \texttt{requestHash}: cùng khóa và cùng nội dung trả lại kết quả cũ; cùng khóa nhưng khác nội dung bị từ chối. Test gửi 10 yêu cầu song song cùng khóa vẫn chỉ tạo một đơn.
\item[Tồn kho đồng thời.] Mỗi món chạy \texttt{UPDATE ... SET stock = stock - n WHERE id = ? AND stock >= n}; nếu không dòng nào được cập nhật thì rollback toàn bộ. Test: 20 người đặt song song khi kho còn 5 thì đúng 5 đơn thành công; hủy đơn hoàn lại tồn kho.
\item[Khuyến mãi và điểm thưởng.] Mã \texttt{WELCOME10} giảm 10\% khi đủ điều kiện. Khi đơn PAID, người dùng được cộng 1 điểm cho mỗi 10.000 đồng; hủy đơn đã PAID thì trừ lại; mỗi lần đổi điểm ghi một bản ghi \texttt{LoyaltyTransaction} (EARN hoặc REVERSAL).
\end{description}

\section{CI}\label{sec:se-ci}
\texttt{ci.yml} chạy cho mọi PR với bốn job: \texttt{api} (cài đặt, lint, unit test, build), \texttt{api-e2e} (migration, seed, test tích hợp trên PostgreSQL thật), \texttt{web} và \texttt{e2e-browser} (Playwright). Lần chạy gần nhất trên \texttt{main} (run \#60) xanh ở cả bốn job. Pipeline CD, triển khai production và quay lui được trình bày trong báo cáo Assignment 2.
